\documentclass[12pt]{article}
\usepackage[utf8]{inputenc}
\usepackage[T1]{fontenc}
\usepackage[margin=1in]{geometry}
\usepackage{setspace}
\usepackage{xcolor}
\usepackage[colorlinks=true, urlcolor=blue!60!black, linkcolor=blue!60!black]{hyperref}
\setlength{\parindent}{0em}
\setlength{\parskip}{0.6em}
\renewcommand{\baselinestretch}{1.0}
\begin{document}
\begin{center}
{\LARGE \textit{Washington v. Glucksberg}, 521 U.S. 702 (1997)\par}
\vspace{0.5em}
{\large Case Brief --- Two-Page Edition\par}
\end{center}

\section{Memory Jogger}
No fundamental right to physician-assisted suicide: the asserted liberty is not ``deeply rooted in this Nation's history and tradition'' (\href{https://supreme.justia.com/cases/federal/us/521/702/}{Justia}).

\section{Detailed Case Facts}
Washington made it a felony to knowingly cause or aid another's suicide attempt, while allowing patients to refuse life-sustaining treatment. Physicians, including Dr. Harold Glucksberg, terminally ill patients, and Compassion in Dying sought a declaration that the ban violated due process.

\section{Procedural History}
The district court and the en banc Ninth Circuit recognized a liberty interest in ``controlling the time and manner of one's death.''

\section{Judicial Votes}
Reversed 9--0. Rehnquist wrote for O'Connor, Scalia, Kennedy, and Thomas; O'Connor concurred; Stevens, Souter, Ginsburg, and Breyer concurred in the judgment.

\section{Holding}
Fundamental rights must be deeply rooted in history and tradition and carefully described. A ``right to commit suicide which itself includes a right to assistance in doing so'' is not fundamental, and the ban rationally serves legitimate interests.

\section{Analysis of Opinions}
\textbf{Majority:} 700 years of Anglo-American law condemned assisting suicide; \textit{Cruzan} rested on the common law of battery; the ban protects life, medical ethics, and vulnerable groups and guards against euthanasia. \textbf{O'Connor:} palliative care was available, so the States' ``laboratory'' should decide. \textbf{Stevens and Breyer:} as-applied claims and a right to die with dignity remain open. \textbf{Souter:} under Harlan's \textit{Poe} method, the interest is serious but outweighed for now.

\section{Examples: Future Applications}
\textbf{Same-Side:} no right to unapproved experimental drugs. \textbf{Opposite-Side:} forcing a ventilator on a competent patient infringes the right recognized in \textit{Cruzan}. \textbf{Fence-Sitter:} bans on surrogacy or embryo disposal, where the level of generality decides.

\section{Critique}
After \textit{Dobbs} used the test to overrule \textit{Roe}, critics say it freezes liberty in the past and clashes with \textit{Lawrence} and \textit{Obergefell}; defenders praise it for leaving end-of-life policy to democratic debate.

\section{Key Quotations}
\begin{quote}
``Our established method of substantive-due-process analysis has two primary features: First, we have regularly observed that the Due Process Clause specially protects those fundamental rights and liberties which are, objectively, `deeply rooted in this Nation's history and tradition,' \ldots\ and `implicit in the concept of ordered liberty,' such that `neither liberty nor justice would exist if they were sacrificed,' \textit{Palko v. Connecticut}, 302 U. S. 319, 325, 326 (1937). Second, we have required in substantive-due-process cases a `careful description' of the asserted fundamental liberty interest.''
\end{quote}
\noindent --- Chief Justice Rehnquist, for the Court, 521 U.S. at 720--721.

\end{document}
